\section{Exact gradient evolution and finite path comparisons}
\label{sec:gradient-path}

We establish two fixed-time approximations for the later path calculus.
The first truncates the exact derivative of a conditional expectation
in operator norm, uniformly in the starting state.  The second replaces
empirical observations along finitely many paths by finite smooth
simulations.  Their error bounds serve different purposes: the first
allows multiplication by a correlated vector, while the second
identifies limiting spatial averages.

We begin with an evolution on all site-vector functions that agrees
with scalar differentiation on gradients.  Its likelihood expansion
controls the first approximation.  Mixed differences then let us
differentiate functions of several stored states.  Finally, density
and entropy comparisons with frozen-field paths give the simulation
approximation.

All intervals in this section are fixed.  Constants may depend on their
lengths, the number of branches, and the finite list of tests, but never
on the dimension or the starting configuration.  The later block
argument will fix these estimates before iterating the intervals.

\subsection{The vector extension and its likelihood expansion}

The derivative of a scalar observable is a site vector.  We first
construct an evolution on arbitrary site-vector functions which agrees
with scalar differentiation on gradients.  For a site vector $b$, write
$D_b=\diag(b_1,\ldots,b_n)$.  The estimates will use the disorder conditions
\begin{equation}
 \begin{gathered}
 \norm J_{\op}\le C_0,\qquad
 \max_i\sum_jJ_{ij}^2\le C_0,\\
 \max_{i,j}|J_{ij}|=o(1),\qquad
 \norm{J\circ J-\beta^2\1\1^t/n}_{\op}=o(1).
 \end{gathered}
 \label{eq:path-disorder}
\end{equation}
Here $\circ$ denotes entrywise multiplication.  These hold with
probability tending to one.  For the last assertion one may truncate
$J_{ij}^2-\beta^2/n$ at $C\log n/n$.  The centered symmetric matrix has
variance parameter $O(1/n)$; the matrix exponential bound gives
$\norm{J\circ J-\E(J\circ J)}=O(\sqrt{\log n/n}+\log^2n/n)$
with probability tending to one.  The discarded probability tends to
zero by Gaussian tails, and the missing diagonal is $O(1/n)$.
The first three assertions follow respectively from the Gaussian matrix
norm bound, chi-square tails for the rows, and a union bound for entries.
The same estimates underlie the independent argument in
\Cref{sec:rcut-introduction}.

\begin{proposition}[Exact vector extension]
\label{prop:path-gradient}
On all vector-valued functions on the cube define
\begin{equation}
 (\mathscr Gp)_i=(L-I)p_i
       -2x_i\sum_j a_{ij}d_jp_i+\sum_j a_{ij}p_j,
 \qquad a_{ij}=d_i m_j,
 \label{eq:path-generator}
\end{equation}
and let $K_t=e^{t\mathscr G}$.  Then $dS_tf=K_tdf$.  Under
\Cref{eq:path-disorder}, uniformly in $x$,
\begin{align}
 a&=JD_v+\frac{2\beta^2}{n}x(mv)^t+o_{\op}(1),
 \label{eq:path-a-expansion}\\
 \norm{aD_v^{-1}}_{\op}&\le C,
 &\max_i\sum_j\frac{a_{ij}^2}{w_j}&\le C.
 \label{eq:path-a-bounds}
\end{align}
The variance factors in the second line are retained even at saturated
fields; no lower bound on $v_j$ or $w_j$ is required.
\end{proposition}

\begin{proof}
\emph{Intertwining.}
The exact product rule is
\[
 d_i(fg)=f\,d_ig+g\,d_if-2x_i(d_if)(d_ig).
\]
Use $Lf=\sum_j(m_j-x_j)d_jf$, $d_i^2f=0$, and
$d_id_jf=d_jd_if$.  Differentiation gives
\[
 d_iLf=Ld_if-d_if+\sum_j a_{ij}d_jf
                         -2x_i\sum_j a_{ij}d_jd_if.
\]
Uniqueness of the finite-dimensional evolution proves intertwining.

\emph{Coefficient bounds with the variance factors retained.}
For $i\ne j$, the subtraction formula for $\tanh$ gives
\[
 \frac{a_{ij}}{v_j}
 =\frac{\tanh(2x_iJ_{ij})}
        {2x_i\{1-m_j\tanh(2x_iJ_{ij})\}}
 =J_{ij}+2x_i m_jJ_{ij}^2+O(|J_{ij}|^3).
\]
The remainder is uniform for $|m_j|\le1$ and
$\max|J_{ij}|\le1/10$, including when the local conditional variance
is small.  Its maximum absolute row and column sums
are $o(1)$, because
$\sum_j|J_{ij}|^3\le\max|J_{ij}|\sum_jJ_{ij}^2$.
The last condition in \Cref{eq:path-disorder} proves
\Cref{eq:path-a-expansion} and the first bound in
\Cref{eq:path-a-bounds}.  Finally,
$|a_{ij}|\le C|J_{ij}|v_j$ and
$v_j=w_j(1+x_jm_j)\le2w_j$ give
$v_j^2/w_j\le2$.  This proves the row bound.
\end{proof}

The weighted coefficient bounds give two controls on this extension:
a pointwise energy inequality for arbitrary vector fields, and a
uniform truncation of its path expansion.  Both will be needed when
conditional derivatives are integrated along a tree.

\begin{proposition}[Rough propagation and path tails]
\label{prop:path-rough}
There are constants $C,c>0$ such that, for every vector function $p$,
\begin{equation}
 (\partial_t-L)\norm{K_tp}^2
 \le C\norm{K_tp}^2-c\sum_jw_j\norm{d_jK_tp}^2.
 \label{eq:path-square}
\end{equation}
In particular, for every starting law $\nu$ and $t\ge0$,
\begin{equation}
 \E_\nu\norm{K_tp}^2\le e^{Ct}\E_{\nu S_t}\norm p^2.
 \label{eq:path-rough-L2}
\end{equation}
The same assertions hold with the multiplication term $ap$ in
\Cref{eq:path-generator} replaced by $z ap$, for each fixed real $z$.

Expand in that multiplication term.  Given $H>0$, there are
$A_H,B_H<\infty$ such that, for $t\ge1$, deleting terms with more than
$A_Ht$ insertions, and then clipping elementary likelihoods at
$\exp(B_Ht)$, changes the resulting map by at most $C_He^{-Ht}$
as a map from $L^2(\nu S_t;\R^n)$ to $L^2(\nu;\R^n)$.
The constants are uniform in $\nu$, including point masses.
Columnwise the same statement holds for Hilbert--Schmidt norms.
For every fixed $t>0$, sending the insertion cutoff and then the
likelihood cap to infinity likewise gives convergence in this same
all-vector-field operator bound.  At $t=0$ the map is the identity.
\end{proposition}

\begin{proof}
We first prove the energy inequality, then represent each row evolution
by a likelihood change, and finally truncate the expansion in the
multiplication term.

\emph{Energy control.}
Put $q=K_tp$ and $\mathcal D(q)=\sum_jw_j\norm{d_jq}^2$.
The jump identity gives
\[
 (\partial_t-L)\norm q^2
 =-2\norm q^2+2q^taq
   -4\sum_{i,j}q_ix_i a_{ij}d_jq_i-2\mathcal D(q).
\]
The middle first-order term in absolute value is at most
$C\norm q\sqrt{\mathcal D(q)}$ by \Cref{eq:path-a-bounds}.
Young's inequality absorbs one copy of $\mathcal D(q)$; the
multiplication term costs $2\norm a\norm q^2$.
Integrating the resulting inequality along an ordinary trajectory
proves \Cref{eq:path-rough-L2}.  Retaining dissipation also gives
\[
 \int_0^t\E_{\nu S_s}\sum_jw_j
          \norm{d_jK_{t-s}p}^2\dd s
 \le C_te^{Ct}\E_{\nu S_t}\norm p^2.
\]

\emph{Likelihood representation and moments.}
To control the expansion tails, represent the evolution without its
multiplication term by a change of the ordinary update rules.  In row $i$
replace, at attempts at $j\ne i$, the field $H_j=(JX)_j$ by
$H_j-2X_iJ_{ij}$; attempts at $i$ are unmodified.  The corresponding
row generator is $L-2x_i\sum_ja_{ij}d_j$.  If $\ell_i[s,t]$ is
its likelihood relative to ordinary attempts, its semigroup, including
$-I$, is
\[
 (T_{t-s}p)_i(x)=e^{-(t-s)}\E_x[\ell_i[s,t]p_i(X_{t-s})].
\]
The site-$i$ spin is part of the ordinary trajectory; only its update
\emph{rule} is left unchanged in this representation.

The required moments come from a one-attempt calculation.  Let
$b_H(\sigma)=e^{\sigma H}/(2\cosh H)$,
$\sigma\in\{-1,1\}$, and $r=b_{H+u}/b_H$.
Taylor's formula and $(\log\cosh)''\le1$ show, for every fixed
real $p$, that
\[
 \E_{b_H}r^p
 =\exp\{\log\cosh(H+pu)-p\log\cosh(H+u)
                    +(p-1)\log\cosh H\}
 \le e^{C_pu^2}.
\]
For a rate-one clock at $j$, the compensator of the $p$-th likelihood
moment is therefore at most $C_pJ_{ij}^2$; the largest entry is small.
The exponential martingale formula gives, conditionally on any past,
\begin{equation}
 \E[\ell_i(I)^p\mid\text{past at the first interval}]
 \le \exp(C_p|I|),
 \label{eq:path-likelihood-moment}
\end{equation}
where $I$ is any fixed union of intervals and $|I|$ is its total
length.  Negative powers are allowed.  All fields and starting states
are covered by the same constants.

\emph{Truncating insertions and then likelihoods.}
Rowwise Cauchy--Schwarz gives
$\norm{T_hp(x)}^2\le e^{Ch}S_h\norm p^2(x)$.
Writing $A$ for multiplication by $a$, the term with $k$ insertions is
\begin{equation}
 \int_{0<t_1<\cdots<t_k<t}
 T_{t_1}A T_{t_2-t_1}A\cdots A T_{t-t_k}p\,
 \dd t_1\cdots\dd t_k.
 \label{eq:path-Duhamel}
\end{equation}
Repeated rowwise Cauchy--Schwarz and the Markov property bound its
$L^2$ norm by $e^{Ct}C^kt^k/k!$ times the terminal $L^2$ norm.
Consequently
\[
 \sum_{k>At}e^{Ct}\frac{(Ct)^k}{k!}
 \le e^{Ct}\left(\frac{eC}{A}\right)^{At}
                   (1-eC/A)^{-1},
\]
which is at most $C_He^{-Ht}$ when $A$ is sufficiently large.
Once this insertion cutoff is fixed, only finitely many likelihood
intervals remain in each term.  For clipping, with any fixed $r>0$,
\[
 \E[\ell_i^2\1_{\ell_i>e^{Bt}}\mid\text{past}]
 \le e^{-rBt}\E[\ell_i^{2+r}\mid\text{past}]
 \le e^{-rBt+C_{2+r}t}.
\]
Apply this estimate at the first clipped interval, use the ordinary
second-moment bounds elsewhere, and sum over at most $At+1$ intervals
and the simplex.  Increasing $B$ after $A$ proves the stated bound.
For fixed $t>0$ and a prescribed error, choose a finite insertion
cutoff making the factorial-series tail sufficiently small.  Then choose
a likelihood cap controlling the finitely many retained terms, using
the displayed moment estimate.  This convergence does not require
$t\ge1$; that restriction was used only in the stated exponential-in-time
tail formulation.  Every estimate sums over columns, which proves the
matrix version.
\end{proof}

\subsection{Mixed differences and conditional derivative matrices}

A multitime calculation is a function of several stored spin states.
Integrating a terminal observation introduces two arguments for its
parent: the already stored state and the starting state of the new
conditional expectation.  If these arguments are denoted by $y,z$,
their identification obeys
\[
 d_j[G(y,z)|_{z=y}]
 =\left[d_j^yG+d_j^zG-2y_jd_j^yd_j^zG\right]_{z=y}.
\]
The last term is a mixed difference.  Separate first-derivative bounds
in each state slot therefore do not close under this operation.

We retain the finite collection of tensors with at most one derivative
in each distinct slot.  Two elementary estimates below control their
products and smooth compositions, after which
\Cref{lem:path-conditional-derivative} carries the collection through
conditioning.  The spatial indices of these tensors are reduced to
path matrices in \Cref{prop:path-opened-derivative}.

Let $F:(\{-1,1\}^n)^m\to\R^n$, with state slots
$x^{(1)},\ldots,x^{(m)}$.  For a nonempty $I\subseteq[m]$ define
\begin{equation}
 (T_IF)_{i,(j_a:a\in I)}
   =\prod_{a\in I}d_{j_a}^{(a)}F_i,
 \qquad C_I(F)=\sup_{x^{(1)},\ldots,x^{(m)}}
                         \norm{T_IF}_{\op}.
 \label{eq:path-mixed-norm}
\end{equation}
All derivative indices are columns of this rectangular matrix; its
only row index is the output site.  Equivalently, $C_I(F)$ bounds the
Euclidean norm of the complete mixed-difference tensor of $c^tF$ by
$C_I(F)\norm c$ for every deterministic output vector $c$.
No bound on the unnormalized Euclidean norm of $F$ is imposed.
Coordinate caps, when needed for labels, are recorded separately.

\begin{lemma}[Products and repeated squares of tensor entries]
Let $B_l$ have $n$ rows and columns indexed by $j_{S_l}$, where the
nonempty sets $S_l$ cover $I$, and suppose $\norm{B_l}_{\op}\le b_l$.\label{lem:path-tensor-product}
Then
\begin{equation}
 \norm{\left(\prod_l B_l(i,j_{S_l})\right)_{i,j_I}}_{\op}
 \le\prod_l b_l.
 \label{eq:path-tensor-product}
\end{equation}
If a rectangular matrix $R$ satisfies
\begin{equation}
 |R_{i,j_I}|
 \le C|B_0(i,j_{S_0})|^2
                     \prod_{l\ge1}|B_l(i,j_{S_l})|,
 \label{eq:path-square-domination}
\end{equation}
with the same covering condition, then
$\norm R_{\op}\le Cb_0^2\prod_{l\ge1}b_l$.
The coefficient implicit in the domination may depend arbitrarily on
all indices.
\end{lemma}

\begin{proof}
\emph{Products.}
First assign independent column indices to every factor.  Their rowwise
tensor product is obtained from $\bigotimes_lB_l$ by restricting its
output to the diagonal copy of $\R^n$, an operator of norm one.
Next require all occurrences of a slot to have the same index.  The
map
\[
 e_{j_I}\longmapsto\bigotimes_l e_{j_{S_l}}
\]
is an isometry, since the sets $S_l$ cover $I$.  This proves
\Cref{eq:path-tensor-product}.  There is no absent, freely repeated
column index that could contribute a dimension factor.

\emph{An entrywise square and arbitrary bounded coefficients.}
For the second assertion put $A=|B_0|^2$ entrywise.  Its maximum row
sum and maximum column sum are at most $b_0^2$, hence the rectangular
Schur test gives $\norm A_{\op}\le b_0^2$.
For any two matrices $A,B$ with the same output rows, their rowwise
tensor product satisfies
\[
 \norm{A\otimes_{\rm row}B}_{\op}
 \le\norm A_{\op}\max_i\norm{\operatorname{row}_i(B)}_2.
\]
For example, factor it as $D_B(A\otimes I)$, where
$(D_Bz)_i=\sum_b B_{ib}z_{i,b}$; the norm of $D_B$ is the displayed
maximum.  Apply this repeatedly to $A$ and the absolute values of the
other factors, whose row Euclidean norms are at most $b_l$.
Restrict repeated column indices by the same isometry as before.
Finally entrywise domination implies
$|u^tRv|\le C|u|^tP|v|$ for the nonnegative dominating matrix $P$.
Taking unit-vector suprema proves the assertion.
\end{proof}

\begin{lemma}[Mixed-slot bounds for finite smooth calculations]
\label{lem:path-mixed-tensors}
Consider a fixed finite calculation whose inputs are finitely many
state slots and fixed site labels.  Its operations are bounded-operator
linear maps on the site index, finite linear combinations, and smooth
componentwise functions with globally bounded derivatives of every
positive order.  Bounded coordinate products may be used through a
fixed smooth extension outside their bounded ranges.  Then every
nonempty tensor in \Cref{eq:path-mixed-norm} has a dimension-independent
bound.  The constants depend on the calculation, the linear operator
bounds, and finitely many scalar derivative bounds.  They are uniform
in fixed past labels whenever those displayed bounds are uniform.
Linear-query values themselves need not be uniformly bounded.
\end{lemma}

\begin{proof}
For spin inputs and fixed labels, the tensors are explicit: a spin
input in slot $a$ has $T_{\{a\}}=I$ and every other nonempty tensor
zero, while a fixed label has all nonempty tensors zero.  Linear maps
and sums preserve the claimed bounds.  The remaining task is closure
under smooth componentwise functions.  We use an exact corner expansion
so that no compact-range assumption is imposed on a linear query.

First consider $F_i=\phi(z_i)$.  Fix $I$, $|I|=r$, and a derivative
tuple $j_I$.  For every nonempty $S\subseteq I$ put
\[
 t_S=(-2)^{|S|}\left(\prod_{a\in S}x_{j_a}^{(a)}\right)
                       (T_Sz)_{i,j_S}.
\]
If the indicated coordinate is flipped in each slot of $B\subseteq I$,
the exact cube identity is
\begin{equation}
 z_i(x^B)=z_i(x)+\sum_{\varnothing\ne S\subseteq B}t_S.
 \label{eq:path-corner-increments}
\end{equation}
Temporarily regard the $2^r-1$ increments as independent scalar
variables.  Repeated use of the fundamental theorem of calculus gives,
for any set $V$ of increments,
\begin{equation}
 \phi\left(z+\sum_{S\in V}t_S\right)
 =\sum_{A\subseteq V}\left(\prod_{S\in A}t_S\right)
   \int_{[0,1]^A}\phi^{(|A|)}
           \left(z+\sum_{S\in A}\theta_St_S\right)\dd\theta.
 \label{eq:path-anchored-expansion}
\end{equation}
The empty term is $\phi(z)$.  Apply this with
$V=\{S:\varnothing\ne S\subseteq B\}$ and take the alternating sum
over $B\subseteq I$ defining $T_I$.  A term survives exactly when
\begin{equation}
 \bigcup_{S\in A}S=I.
 \label{eq:path-covering}
\end{equation}
Indeed, otherwise summing over a missing slot cancels the term.
The remaining normalization is $2^{-r}$ and a product of column signs.

Every surviving monomial thus contains enough increments to cover all
derivative slots.  Its coefficient may still depend on the derivative
indices, so the product estimate alone is not yet sufficient.  Expand
that integral coefficient once at $z$:
\[
 \int\phi^{(|A|)}\left(z+\sum_{S\in A}\theta_St_S\right)\dd\theta
 =\phi^{(|A|)}(z)+\sum_{S\in A}t_SR_{A,S}(z,t),
 \qquad |R_{A,S}|\le\norm{\phi^{(|A|+1)}}_\infty.
\]
This is an exact integral remainder.  The first term is a bounded
row coefficient times a product of tensors satisfying
\Cref{eq:path-covering}, so \Cref{eq:path-tensor-product} applies.
The remainder acquires a second copy of an increment already present
in its monomial.  It is therefore dominated by a product with one
squared tensor factor,
and \Cref{eq:path-square-domination} applies, even though its coefficient
depends on every derivative index.  There are finitely many terms.

For a componentwise function of $q$ scalar inputs $z^1,\ldots,z^q$,
introduce increments $t_{b,S}$ for every channel $b$ and nonempty
$S\subseteq I$.  Apply the same anchored expansion in these
$q(2^r-1)$ formal variables.  The surviving monomials again cover
$I$, and expanding each coefficient once repeats an increment already
in that monomial.  Bounded mixed scalar derivatives through order
$q(2^r-1)+1$ suffice at this gate.  The same two estimates prove the
claim.  Induction over the fixed calculation completes the proof.
\end{proof}

\begin{lemma}[Conditional derivative matrices]
\label{lem:path-conditional-derivative}
Let $F$ be a finite multitime site-vector calculation satisfying the
mixed-slot bounds of \Cref{lem:path-mixed-tensors}.  Evaluate it on a
fixed finite ordinary path tree, condition at one joining state $x$,
and hold the realized past fixed.  Write
$A(x)=d_x\E_xF$, with output as row and derivative index as column.
There are common finite Duhamel and likelihood truncations $A_M$ such
that
\begin{equation}
 \sup_x\norm{A(x)-A_M(x)}_{\op}\longrightarrow0.
 \label{eq:path-conditional-op}
\end{equation}
The assertion is uniform over stored parameters with uniform mixed-slot
bounds.  In particular, for any vector $u$, possibly depending on the
joining state and stored past,
\begin{equation}
 \norm{D_u(A-A_M)}_{\HS}
 \le\norm{A-A_M}_{\op}\norm u.
 \label{eq:path-row-weighted-error}
\end{equation}
The same statement holds on an added stationary reverse branch.
Actual and predictor contributions may use identical subdivisions and
common truncations, and are recombined before a centering test.
\end{lemma}

\begin{proof}
We integrate one leaf at a time.  The propagation bound controls the
evolving derivative slot, while the mixed tensors retain the derivatives
in every stored slot.  After writing this as a finite linear recursion,
we use one common truncation of each evolution operator.

\emph{Propagation with auxiliary derivative indices.}
The point-start bounds in \Cref{prop:path-rough} are
\begin{align}
 \norm{K_hq(x)}&\le B_h\{S_h\norm q^2(x)\}^{1/2},\notag\\
 \norm{(K_h-K_{h,M})q(x)}
   &\le\varepsilon_{h,M}\{S_h\norm q^2(x)\}^{1/2},
 \qquad \varepsilon_{h,M}\longrightarrow0.
 \label{eq:path-amplified-tail}
\end{align}
They amplify to any auxiliary Euclidean space by applying them to
each component and summing squares.  In particular the other derivative
indices of $T_IF$ can all be retained as Hilbert coordinates without
a dimension factor.  The truncation is one common linear operator,
not a choice depending on an output test vector.

\emph{Integrating a leaf and identifying its parent.}
Integrate a final observed slot $m$ from its parent slot $p$ over an
interval of length $h$.  Let $y$ denote all slots other than $m$.
Before identifying the parent and the evolving starting state, set
\[
 G(y,z)=S_h[w\mapsto F(y,w)](z),
 \qquad H(y)=G(y,z)\big|_{z=y^{(p)}}.
\]
The two occurrences of the parent are independent arguments of $G$.
The copy identity at the start of this subsection gives
\begin{equation}
 d_j^{(p)}H
 =\left[d_j^{(p)}G+d_j^zG
           -2y_j^{(p)}d_j^{(p)}d_j^zG\right]_{z=y^{(p)}}.
 \label{eq:path-copy-difference}
\end{equation}
All other frozen-slot differences commute with $S_h$ and $K_h$.
For $p\notin I$, the resulting tensor is simply $S_hT_IF$ restricted
to $z=y^{(p)}$.  If $p\in I$ and $I'=I\setminus\{p\}$, its three
terms are respectively
\begin{equation}
 S_hT_IF,\qquad
 K_hT_{I'\cup\{m\}}F,\qquad
 -2D_{y^{(p)}}\,\mathsf R_{j_m=j_p}
                         K_hT_{I\cup\{m\}}F.
 \label{eq:path-copy-tensors}
\end{equation}
The sign multiplier in the last term acts on the derivative index.
In the last two terms $K_h$ acts only on the slot-$m$ derivative
coordinate; all other indices are auxiliary Hilbert coordinates.
The restriction $\mathsf R_{j_m=j_p}$ is a coordinate restriction
within that Hilbert space, hence has norm one.  It never identifies
an output row with an input column across the chosen flattening.

We can now estimate the three terms without changing the tensor
flattening.  Fix a deterministic output vector $c$ and apply the
transpose of each tensor to $c$.  Jensen for $S_h$ and
\Cref{eq:path-amplified-tail} give
\begin{equation}
 C_I(H)\le C_I(F)+B_hC_{I'\cup\{m\}}(F)
                          +2B_hC_{I\cup\{m\}}(F)
 \quad(p\in I),
 \label{eq:path-copy-bound}
\end{equation}
and $C_I(H)\le C_I(F)$ otherwise.  Taking the supremum over $c$
after these deterministic inequalities proves the operator bounds.
Only one difference per original slot occurs.  Iterate this finite
elimination along a path, or remove leaves of a finite causal tree.
Other branches can be held as parameters during each leaf integration.
If the original calculation does not use the joining state, adjoin it
as a dummy slot.  Its initial difference tensors are zero.  Reversibility
makes a stationary reverse edge another use of the same $S_h$.

\emph{One truncation for the complete derivative matrix.}
The leaf elimination is a finite linear recursion on the family of mixed
tensors.  Open that recursion and replace each
$K_h$ by the same chosen $K_{h,M}$ at every occurrence with that
interval.  Ordinary $S_h$ expectations may remain exact.  Each step
has the bounded norm in \Cref{eq:path-copy-bound}; replacing one step
costs at most the corresponding $\varepsilon_{h,M}$ from
\Cref{eq:path-amplified-tail}.  Telescoping the finitely many steps
gives \Cref{eq:path-conditional-op}.  This approximates the derivative
of the exact conditional expectation.  It does not assume that $A_M$
is the derivative of a separately approximated scalar message, and
no operator supremum is moved through an expectation.

\emph{Weighted errors and recombination.}
Finally,
$\norm{D_uE}_{\HS}^2=\sum_i|u_i|^2\norm{\operatorname{row}_i(E)}^2$
proves \Cref{eq:path-row-weighted-error} pointwise.  The recursion
and the common truncation are linear in the endpoint observables;
therefore actual and predictor terms can be recombined before the
centered label is used.  Opening their ordinary conditional expectations
as causal branches is done only after these uniform norm errors have
been removed.
\end{proof}

\subsection{Admissible trees and the order of limits}

We have controlled conditional derivatives pointwise.  To compare the
path observables themselves, we must specify which conditional path
laws and weights are allowed.  The following definition keeps each
new branch attached at a single endpoint.

\begin{definition}[Finite tree tests]
\label{def:path-tree}
A stationary path tree has finitely many edges of fixed finite lengths.
Choose a root state with law $\mu$ and run independent ordinary
heat-bath paths outwards along the edges.  Reversibility permits an
edge to be read in either direction.  A nonstationary tree is obtained
from a forward path whose initial density $h$ satisfies
$\log\norm h_\infty=o(n)$, and by making finitely many conditional
copies of portions at a shared \emph{single endpoint}, using their
actual conditional laws.  No new branch is conditioned on both of its
endpoints.

A test has a fixed finite list of times and a fixed polynomial bound
$n^M$.  Its weight $W_n$ satisfies $|W_n|\le n^M$ and
$\E|W_n|\le C$.  A bounded test assertion is called typical if it holds
in expectation with every such weight, on disorder events of probability
tending to one.  One may pass to subsequences and successive finite
approximations.  The exponent $M$, all times, the tree, and each
approximation are fixed before $n\to\infty$.
\end{definition}

The definition also covers weighted squared norms: the weight can be
$\norm u^2$ when $u$ is polynomially bounded with bounded mean square.
Statements involving integrated auxiliary times mean integration over a
fixed finite-measure parameter set first.  They do not permit choosing
a worst parameter from an exceptional set of measure zero.

\begin{lemma}[Transfer to actual conditional trees]
\label{lem:path-tree-transfer}
A stationary tree event of probability $e^{-cn}$ has negligible weighted
probability in \Cref{def:path-tree}.  The same holds for a finite family
of such events and for exponentially small integrated bad mass in a
fixed finite-measure parameter space.
\end{lemma}

\begin{proof}
Forward evolution preserves the density bound.  Reverse conditioning
introduces denominators, so we first remove joining states where those
denominators are too small.  All forward marginals have relative density
at most $e^{o(n)}$, since
$S_t$ is a positive $L^\infty(\mu)$ contraction.  A reverse conditional
edge from time $t$ to $s$ has, relative to the stationary reverse edge,
the factor $h_s(X_s)/h_t(X_t)$.  Discard endpoints where
$h_t<e^{-\delta n}$.  Their actual marginal probability is at most
$e^{-\delta n}$, and their polynomially weighted probability vanishes.
The remaining joint density on a tree with $r$ reverse edges is at most
$\exp\{o(n)+r\delta n\}$.  Choose $\delta<c/(2r+2)$ and compare
probabilities.  The same calculation after Fubini proves the integrated
assertion.  For a finite tree with more than one original marginal,
apply this argument successively at each copying operation.
\end{proof}

\subsection{Two comparisons with frozen fields}

Freezing the field over each mesh bin makes its conditional endpoint
law a product measure.  We will need both a bound on the log density
relative to this reference and a sharper entropy bound for empirical
tests.  Choose a mesh containing all observation and branching times.
On a mesh bin
$[s,s+h]$, its frozen-field law keeps the fields $JX_s$ fixed while
independent rate-one clocks refresh the sites.  Conditional on $X_s$,
the endpoint law is a product measure: site $i$ remains at its old
value with probability $e^{-h}$ and otherwise has law $b_{(JX_s)_i}$.
In particular its probability of changing is at most $h$.

All reference laws below are conditional on the fixed disorder.
Let $P$ be the stationary ordinary marked-attempt law on the chosen
root-directed tree: its root has law $\mu$, and child paths are
conditionally independent ordinary evolutions.  Let $P_h$ have the
same root law and tree factorization, with fields frozen separately
in each mesh bin.  Stationary reversal identifies the ordinary
reference on a reversed edge; an actual nonstationary reverse edge
still uses its actual conditional law.  For a bin $b$ of length
$\delta_b$, the endpoint kernel of $P_h$ is the product measure
\begin{equation}
 \pi_b(x,\dd y)=\bigotimes_{i=1}^n
 \left[e^{-\delta_b}\delta_{x_i}(\dd y_i)
   +(1-e^{-\delta_b})b_{(Jx)_i}(\dd y_i)\right].
 \label{eq:path-frozen-product-law}
\end{equation}
Thus $P_h$ is a tree path law and $\pi_b(x,\cdot)$ is one of its
conditional endpoint laws.  Each coordinate of the latter has minority
mass at most $\delta_b\le h$.

Write $Q$ for a retained, normalized admissible actual tree law,
including a permitted nonnegative polynomial weight and conditioning
on an event of polynomially significant probability.  For a bounded
test, a weighted part with total mass below $n^{-1}$ is negligible
before normalization; all other normalizers are at least inverse
polynomial.  Signed weights are split into positive and negative
parts.  All superpolynomial discards are made on the original actual
forward marginals before forming $Q$.

\begin{lemma}[Density and entropy comparisons]
\label{lem:path-two-mesh}
Fix a time horizon and a finite tree.  For sufficiently small fixed
mesh size $h$, the following comparisons hold.

First, after events of probability smaller than every fixed inverse
power of $n$ have been discarded on each actual forward marginal, the
log density between actual and frozen paths, read in the tree-rooted
direction, is at most $\kappa(h)n+o(n)$ in absolute value, where
$\kappa(h)\downarrow0$.

Second, for spatial empirical tests one has the stronger entropy
estimate
\begin{equation}
 H(Q\mid P_h)\le C_T hn+o(n),
 \label{eq:path-mesh-entropy}
\end{equation}
after the prescribed count and field truncations.  Here $P_h$ is the frozen tree law.  The law $Q$ is an admissible actual
tree law, with the polynomial weights of \Cref{def:path-tree}, possibly
conditioned on an event of probability at least $n^{-M}$ for fixed $M$.
The superpolynomial discards are made first on its original forward
marginals.  On the retained set its additional entropy relative to the
stationary comparison is $o(n)$.  Truncation errors vanish for this
class in the successive order $n\to\infty$, cutoff $\to\infty$ at
fixed $h$, then $h\downarrow0$.
For use at a fixed $h$ the cutoff is chosen so that its entropy error
is at most $hn$.
\end{lemma}

\begin{proof}
We first compare absolute log densities on retained forward and reverse
bins.  We then retain the quadratic log-likelihood cost to obtain the
sharper entropy estimate.  In both arguments the exceptional events are
discarded on the original forward marginals before forming the weighted
conditional law.

\emph{The forward log-density bound.}
Write $\Delta(u)=J(X_u-X_s)$ in a bin starting at $s$, and let
$N_{i,b}$ be its attempt counts.  On a fixed mesh, Poisson exponential
bounds imply, with exponential accuracy,
\begin{equation}
 \sup_{u\in b}\norm{X_u-X_s}^2\le C hn,
 \qquad
 \sup_{u\in b}\norm{\Delta(u)}^2\le C hn.
 \label{eq:path-bin-square}
\end{equation}
The log ratio at an attempt at $i$ is bounded by $2|\Delta_i|$.
For a bounded predictable integrand the compensated attempt sum obeys
\[
 \P(M\ge a)\le
 \exp\!\left(-\frac{a^2}{2(V+ba/3)}\right),
\]
where $V$ bounds its predictable bracket and $b$ its jumps.  Apply this
after cutting $|\Delta_i|$ at $B$, then use
$\sum_i|\Delta_i|\le\sqrt n\norm\Delta$ for its compensator.
The bounded part costs $O_T(\sqrt h)n$.
For the remaining part, $\sum_i\Delta_i^2\le Chn$ bounds the
compensator by $Chn/B$.  Its individual jumps are $O(\sqrt n)$,
its bracket is $O_T(n)$, and the displayed inequality makes a fixed
positive multiple of $n$ an event of superpolynomially small probability.
This proves the forward density comparison.

\emph{Reversed bins.}
For a reversed bin separate the forward within-bin difference from
$J(X_{s+h}-X_s)$.  The latter vector has norm $O(\sqrt{hn})$.
Its entries larger than $B$ occupy at most $Chn/B^2$ sites.  Independently
of how this set is selected, the sum of $N_{i,b}^2$ over such a set is
bounded by the sum of the largest $Chn/B^2$ squared Poisson counts.
For each fixed $h$, truncating counts at a level $K$ gives
\[
 \sum_{i\in I}N_{i,b}^2
 \le K^2|I|+\sum_iN_{i,b}^2\1_{N_{i,b}>K}
 =o_B(hn)
\]
with superpolynomial accuracy, by first choosing $K$ and then $B$.
Cauchy--Schwarz bounds the across-bin excess by $o_B(hn)$.
The bounded part has exponential stationary tails by reversibility.
Discard the excess on each actual edge in its original forward
direction, before applying \Cref{lem:path-tree-transfer} to the bounded
part.  Summing the finitely many bins proves the first assertion.

\emph{The entropy of the retained actual tree.}
For the entropy assertion, first compare $Q$ with $P$.  Let
$c_n=o(n)$ bound the logarithms of the original forward marginal
densities.  Discard reverse joining states with marginal density below
$e^{-a_n}$, where $a_n=o(n)$ and $a_n/\log n\to\infty$.
Their actual marginal probabilities are at most $e^{-a_n}$ and remain
negligible under the permitted polynomial weights.  On the retained
set, the reverse density factors $h_s(X_s)/h_t(X_t)$ bound the total
tree log density by $C(c_n+a_n)$.  Polynomial reweighting and retained
event normalization add $O(\log n)$.  Thus $H(Q\mid P)=o(n)$.
The change of reference is exactly
\begin{equation}
 H(Q\mid P_h)=H(Q\mid P)
                   +\E_Q\log\frac{\dd P}{\dd P_h}.
 \label{eq:path-entropy-reference}
\end{equation}
\emph{Changing to the frozen reference.}
The root laws cancel in the last ratio, leaving the root-directed
within-bin refresh log ratios.  The absolute estimate already proved
would lose the desired factor $h$.  We instead keep the quadratic
log-cosh remainder.  Conditional on the old
configuration, the Kullback--Leibler divergence
between two refresh coins at fields $H$ and $H+u$ is at most $u^2/2$.
Consequently its predictable integral over all bins is at most
\[
 C\int_0^T\norm{J(X_u-X_{\lfloor u\rfloor_h})}^2\dd u
 \le C_Thn
\]
on \Cref{eq:path-bin-square}.  For the centered log likelihood $M$ with field differences cut at $B$,
stop when \Cref{eq:path-bin-square} fails.  Its jumps are bounded by
$CB$ and its predictable bracket is at most $C_T hn$: the conditional
variance of a refresh log ratio is at most a constant times the squared
field difference.  The bounded-jump exponential martingale inequality
therefore gives
\[
 \log\E_P e^{\theta M}\le C_T\theta^2 hn,
 \qquad |\theta|B\le c.
\]
Apply the entropy inequality
$\E_QM\le\theta^{-1}\{H(Q\mid P)+\log\E_Pe^{\theta M}\}$
with fixed positive $\theta\le c/(1+B)$.  At fixed $h,B$, the
additional $o(n)$ entropy costs $o(n)$, and the centered term costs
$O_T(hn)$.  The same estimate for $-M$ gives the reverse bound.
It remains to remove the field cutoff.  For the unbounded excess use
the preceding forward compensator bound and the reversed-bin count bound.
At fixed $h$ its normalized entropy
contribution tends to zero with the cutoff; choose it below $h$ before
taking the next mesh limit.  Conditioning on any retained event of
polynomially significant probability adds $O(\log n)$ entropy.
All superpolynomial discards remain negligible under the permitted
polynomial weights and polynomially significant conditioning.  An
arbitrary $o(n)$-entropy tilt would not have this property and is not
being allowed here.  The chain rule over the finite tree, with the same
argument on every root-directed edge, proves \Cref{eq:path-mesh-entropy}.
\end{proof}

The absolute-density estimate is sufficient for subsequent Gaussian
operator comparisons.  It is not sufficient for refining empirical
path tests.  The extra factor $h$ in \Cref{eq:path-mesh-entropy} supplies
that refinement through the next elementary product estimate.

\begin{lemma}[Rare-update product concentration]
\label{lem:path-product}
Let $\pi$ be a product measure on $\{-1,1\}^n$ such that each
coordinate has minority probability at most $h<1/4$.  If
$\sup_x\norm{dF(x)}\le B/\sqrt n$, then
\begin{equation}
 \log\E_\pi e^{\theta(F-\E_\pi F)}
 \le\frac{CB^2\theta^2}{n\log(1/h)}.
 \label{eq:path-product-mgf}
\end{equation}
Consequently, for every $Q\ll\pi$,
\begin{equation}
 |\E_QF-\E_\pi F|
 \le\frac{CB}{\sqrt{n\log(1/h)}}\sqrt{H(Q\mid\pi)}.
 \label{eq:path-product-entropy}
\end{equation}
\end{lemma}

\begin{proof}
We prove a two-point entropy estimate and then tensorize it.  The
small minority mass enters through the factor $\log(1/p)$; keeping
that factor is what makes the subsequent mesh refinement converge.

\emph{One coordinate.}
For a nondegenerate Bernoulli coordinate let its minority mass be
$p\in(0,1/4]$, put $L_p=\log(1/p)$, and set
$\Lambda_p(t)=\log(1-p+pe^t)-pt$.  For $|t|\le L_p/2$,
\[
 \Lambda_p''(t)
 =\frac{p(1-p)e^t}{(1-p+pe^t)^2}
 \le C\sqrt p\le C/L_p.
\]
Since $\Lambda_p(0)=\Lambda_p'(0)=0$, Taylor's formula gives
$\Lambda_p(t)\le Ct^2/L_p$ in this range.  Outside it,
$\Lambda_p(t)\le|t|\le2t^2/L_p$.  Jensen also gives
$\Lambda_p(t)\ge0$.

If a two-point function $G$ has values $a,a+t$, with the second value
on the minority atom, direct calculation yields
\[
 \frac{\Ent_p(e^G)}{\E_p e^G}
       =t\Lambda_p'(t)-\Lambda_p(t).
\]
Convexity at $t$, evaluated at $2t$, gives
$t\Lambda_p'(t)\le\Lambda_p(2t)-\Lambda_p(t)$ for either sign of $t$.
Consequently
\begin{equation}
 \frac{\Ent_p(e^G)}{\E_pe^G}
 \le\Lambda_p(2t)-2\Lambda_p(t)
 \le\frac{Ct^2}{\log(1/p)}.
 \label{eq:path-bernoulli-entropy}
\end{equation}
This obtains the entropy inequality by convexity, without
differentiating a bound on $\Lambda_p$.  A deterministic coordinate
contributes zero entropy.

\emph{Tensorization and the exponential-moment bound.}
Apply \Cref{eq:path-bernoulli-entropy} conditionally with
$G=\theta F$ and $|t|=2|\theta d_iF|$.  Entropy tensorization gives
\[
 \Ent_\pi(e^{\theta F})
 \le\frac{C\theta^2}{\log(1/h)}
       \E_\pi\left[e^{\theta F}\sum_i(d_iF)^2\right]
 \le\frac{CB^2\theta^2}{n\log(1/h)}\E_\pi e^{\theta F}.
\]
For $\psi(\theta)=\log\E_\pi e^{\theta F}$ this bounds
$\theta\psi'(\theta)-\psi(\theta)$.  Integrate the derivative of
$\psi(\theta)/\theta$ from zero to obtain
\Cref{eq:path-product-mgf}; use $-F$ for negative $\theta$.

\emph{Changing the product law.}
Finally, the entropy variational inequality gives
\[
 \theta(\E_QF-\E_\pi F)
 \le H(Q\mid\pi)+\frac{CB^2\theta^2}{n\log(1/h)}.
\]
Optimize over $\theta>0$ and apply the same argument to $-F$.
This proves \Cref{eq:path-product-entropy}.
\end{proof}

\begin{proposition}[Finite simulations and endpoint predictors]
\label{prop:path-simulation}
Every bounded smooth empirical test of finitely many states on an
admissible tree has the same limiting value as successively finer,
finite smooth simulations from the corresponding endpoint states.
For a single observation interval and a test with difference norm
$B/\sqrt n$, the error is at most
\begin{equation}
 \frac{C_TB}{\sqrt{\log(1/h)}}+o_n(1).
 \label{eq:path-simulation-error}
\end{equation}
The statement holds for replicated conditional tests, using the same
predictor approximation on each copy.  A linear query of bounded smooth
outputs may be included after clipping; its spatial square tails are
uniformly negligible in the successive approximation limits.
\end{proposition}

\begin{proof}
We first compare one terminal test under the actual and autonomous
frozen laws.  This requires two telescopes with a common ordinary
semigroup reference.  We then treat stored states by backward leaf
elimination and smooth each fixed simulation.

\emph{The actual-law telescope.}
Choose a regular $h$-mesh augmented by the finitely many prescribed
observation times, so that $0=t_0<\cdots<t_N=T$,
$\delta_b=t_b-t_{b-1}\le h$, and $N\le C_T/h$.
First consider one terminal test $F$, with
$\sup_x\norm{dF(x)}\le B/\sqrt n$, and put
\[
 f_b=S_{T-t_b}F,\qquad
 \mathsf P_b(x,\cdot)=S_{\delta_b}(x,\cdot).
\]
The point-start bound in \Cref{prop:path-rough} gives
$\sup_x\norm{df_b(x)}\le C_TB/\sqrt n$.
Project $Q$ and $P_h$ to the mesh states and put
$q_b(\cdot\mid\mathcal F_{b-1})
 =Q(X_{t_b}\in\cdot\mid\mathcal F_{b-1})$, where
$\mathcal F_{b-1}=\sigma(X_{t_0},\ldots,X_{t_{b-1}})$.
The entropy chain rule and data processing give
\[
 \sum_{b=1}^N\E_Q
 H\!\left(q_b(\cdot\mid\mathcal F_{b-1})
       \mid\pi_b(X_{t_{b-1}},\cdot)\right)
 \le H(Q\mid P_h)\le C_Thn+o(n).
\]
The omitted initial-state entropy is nonnegative.
Since $f_{b-1}=\mathsf P_bf_b$, the exact telescope is
\begin{equation}
 \E_QF(X_T)-\E_Q S_TF(X_0)
 =\sum_{b=1}^N\E_Q[(q_b-\mathsf P_b)f_b].
 \label{eq:path-actual-telescope}
\end{equation}
Both kernels on the right are evaluated at the actual joining
history; $\mathsf P_b$ depends on it only through $X_{t_{b-1}}$.
Insert $\pi_b$ between them and apply
\Cref{lem:path-product} separately to $q_b$ and $\mathsf P_b$.
The second entropy has the uniform bound
\begin{align*}
 H(\mathsf P_b(x,\cdot)\mid\pi_b(x,\cdot))
 &\le H(\text{ordinary bin path from }x
                  \mid\text{frozen bin path from }x)\\
 &\le C\int_0^{\delta_b}\E_x\norm{J(X_u-x)}^2\dd u
 \le Cn\delta_b^2.
\end{align*}
The first inequality is data processing, the second is the
refresh-coin log-cosh bound, and the last uses
$\E_x\norm{X_u-x}^2\le4nu$ and the bound on $\norm J$.
In particular it holds for every joining-state law.
Writing $H_b$ for the sum of these two conditional entropies yields
\begin{align*}
 |\E_QF(X_T)-\E_Q S_TF(X_0)|
 &\le\frac{C_TB}{\sqrt{n\log(1/h)}}
       \sum_{b=1}^N\E_Q\sqrt{H_b}\\
 &\le\frac{C_TB}{\sqrt{n\log(1/h)}}
       \sqrt{N\sum_{b=1}^N\E_QH_b}
 \le\frac{C_TB}{\sqrt{\log(1/h)}}+o_n(1).
\end{align*}
Here dimension tends to infinity at fixed $h$.

\emph{The autonomous frozen-law telescope.}
Let $\widehat P_h$ be the autonomous frozen simulation started from
the same initial law $Q_0$.  The second, zero-tilt telescope is
\begin{equation}
 \E_{\widehat P_h}F(X_T)-\E_{Q_0}S_TF(X_0)
 =\sum_{b=1}^N\E_{\widehat P_h}
                         [(\pi_b-\mathsf P_b)f_b].
 \label{eq:path-frozen-telescope}
\end{equation}
Its conditional kernels equal $\pi_b$, so their entropy relative
to $\pi_b$ is zero.  The ordinary-kernel entropies still sum to at
most $Cn\sum_b\delta_b^2\le C_Thn$, uniformly in $Q_0$.
Consequently this second difference is at most
$C_TB/\sqrt{\log(1/h)}$.
Since $\E_QS_TF(X_0)=\E_{Q_0}S_TF(X_0)$, the triangle
inequality proves \Cref{eq:path-simulation-error}, after enlarging
$C_T$, for the difference between the actual and frozen tests.

\emph{Stored states and admissible conditioning.}
Stored parameters may be included by writing the test as $F(y,x)$
with the same derivative bound uniformly in $y$.
Disintegrate the reference tree at the joining state and the passive
variables $y$.  The active future must retain its ordinary,
respectively frozen, Markov kernels under that reference law.
Past variables and other branches conditionally independent at the
joining state have this property.  Conditioning on a later endpoint
of the active branch generally does not.  The conditional entropy
of the active future is bounded by the total tree entropy by the
chain rule, so both telescopes apply after averaging over $(y,X_0)$,
and the simulation retains their joint initial law.
A backwards leaf elimination gives precisely these admissible
disintegrations.  For replicated predictions the other copies are
passive and the test bound is uniform in their realized values.

\emph{Smoothing a fixed simulation.}
A fixed frozen simulation is a finite calculation using $J$, independent
site clocks and uniforms, and scalar threshold maps.  Replace each
threshold by a smooth transition in a strip of width $\varepsilon$.
Fresh uniforms put at most $O(\varepsilon)$ mass in that strip.
Induction over finitely many operations, using the bound on $\norm J$,
shows that the normalized mean square error tends to zero.  Derivative
constants of the resulting simulation can depend on its fixed mesh and
smoothing width.

\emph{Backward replacement and derivative control.}
Replace the last observed state first.  The smooth simulation produces
a new smooth test of earlier states.  Fix that calculation and its
finite derivative constants before choosing the accuracy at the next
earlier replacement.  At a branch, duplicate the branching state as
an input to the child calculations.  For replicated predictions, hold
all other copies as passive parameters while replacing one copy; the
conditional entropy chain rule bounds its cost.  Thus its derivative
constant depends only on its own currently fixed test, not on the
accuracy subsequently selected for a passive copy.  This backwards,
nested construction avoids requiring constants uniform over all meshes.

For clarity, the derivative bound for a fixed smooth recipe follows
without any probabilistic assertion.  If a vector $z$ has bounded
half-difference matrix, then, off its diagonal,
\[
 d_j\phi(z_i)=\phi'(z_i)d_jz_i
                  +O(|d_jz_i|^2).
\]
The absolute row and column sums of the error are bounded by the
squared operator norm of that matrix.  Keep its diagonal exact.
The Schur estimate in \Cref{lem:path-tensor-product} and
induction give a bounded derivative matrix.  A normalized empirical
inner product has difference norm $O(n^{-1/2})$ by the product rule.
The remaining assertion concerns square tails of linear queries.
Test such a tail with a smooth excess
function whose second derivative is bounded and whose first derivative
is at most $C(1+|z|)$.  The same Taylor calculation, using bounded column
square sums, gives a difference bound independent of the tail cutoff.
For each fixed simulation Gaussian regression gives square integrability
and vanishing square tails.  The uniform comparison just proved then
transfers this assertion to the original paths.
\end{proof}
